\centerline{\bf \Huge Заряженый листочек}

\begin{flushright}
        \scriptsize{ссылка на \href{https://turgor.ru/lktg/2019/5-Discharging%20method/index.htm}{проект ЛКТГ}}
\end{flushright}

\problem{1}
В прямоугольной таблице $m$ строк и $n$ столбцов ( $m<n$ ). В некоторых клетках таблицы стоят звездочки, так что в каждом столбце стоит хотя бы одна звёздочка. Докажите, что существует хотя бы одна такая звездочка, что в одной строке с нею находится больше звездочек, чем с нею в одном столбце.

\problem{2}
В ряд стоят $n$ коробок. В них суммарно лежит $n$ шариков. За один ход можно выбрать коробку и переложить из нее в каждую соседнюю коробку по одному шарику. Докажите, что, какие бы операции ни делались, рано или поздно хотя бы один шарик окажется в крайней справа коробке.

\problem{3} 
По окружности отметили 40 красных, 30 синих и 20 зеленых точек. На каждой дуге между соседними красной и синей точками поставили цифру 1 , на каждой дуге между соседними красной и зеленой - цифру 2 , а на каждой дуге между соседними синей и зеленой - цифру 3. (На дугах между одноцветными точками поставили 0 .) Найдите максимальную возможную сумму поставленных чисел.

\problem{4}
В некоторых клетках клетчатой полосы $1 \times n$ стоят фишки. В первой клетке (самой левой) $x_1$ штук, во второй $-x_2, \ldots$, в $n$-й клетке $-x_n$. Двое игроков играют в следующую игру: каждым ходом первый игрок выбирает некоторое подмножество фишек, а второй либо передвигает все фишки этого подмножества на 1 клетку влево, а остальные убирает с доски, либо делает то же самое с дополнением выбранного подмножества. Первый хочет, чтобы какая-то фишка оказалась в самой первой клетке, а второй хочет ему помешать. При каких $x_1, x_2, \ldots, x_n$ у первого игрока есть выигрышная стратегия?

\problem{5}
В квадрат $10 \times 10$ по порядку расставлены числа от 1 до 100. За ход можно выбрать число, и прибавить к нему 2 , а из двух соседей, лежащих с ней на одной прямой, вычесть по 1 . Или, наоборот: отнять 2 и прибавить к двум соседям по 1 . В конце снова получилась расстановка чисел от 1 до 100 . Докажите, что она совпала с исходной.


\newpage

\hspace{3mm}

\problem{6}
На бесконечной в обе стороны полосе из клеток, пронумерованных целыми числами, лежит несколько камней (возможно, по нескольку в одной клетке). Разрешается выполнять следующие действия:

--- Снять по одному камню с клеток $n-1$ и $n$, и положить один камень в клетку $n+1$;

--- Снять два камня с клетки $n$ и положить по одному камню в клетки $n+1, n-2$.

(a) Докажите, что при любой последовательности действий мы достигнем ситуации, когда указанные действия больше выполнять нельзя;

(b) Докажите, что конечная ситуация не зависит от последовательности действий (а зависит только от начальной раскладки камней по клеткам).

\problem{7}
На плоскости дано $n$ окружностей радиуса 1 , причем известно, что каждая пересекается хотя бы с одной другой окружностью, и никакая пара не касается. Докажите, что все вместе окружности образуют не меньше $n$ точек пересечения (в одной точке могут пересекаться более двух окружностей).